% Pursuing Stacks (cotes 134-2 to 134-8): archive cote 134-7, site page 36 (PDF page 37), typescript page 478.
% typescript. Draft compared: ps6.tex:1723-1793.
% AI-transcribed by Claude Opus 5.5 (claude-opus-5-5, adaptive thinking, effort high), Message Batches API, batch ps3 (msgbatch_01XdFmuVVBAcgc7EFxdtCqsj), 2026-11-12; not yet checked by eye.
\scanpage{134-7/36}
\mspage{478}

where the L in LH. suggests that we are taking something similar to a
total left derived functor, and where definitely in the right hand mem-%
ber we had to write L$_{\mathrm{pt}}$ and not L$^{\bullet}$ , in order not to \xout{\unsure{sink into}} get sunk
into a morass of confusion ! In the formul\unsure{ae} (U) we should assume \over{\unsure{n}}{h}ow-
\unsure{ever} that X$_{*}$ is a smooth Kan complex, which will imply (if indeed \xout{H}
Hot$_n$(k) may be described in terms of sK\hand{$\mathcal{M}$}$_n$(k) as said above) \unsure{that} LH.
may be viewed as a functor
\[ \text{(U')} \qquad \mathrm{Hot}_n(k) \longrightarrow \mathrm{D.}(\mathrm{Ab}_k) \quad , \]
and likewise the H$_i$'s are functors from Hot$_n$(k) to Ab$_k$ . In order to
compute these homology invariants for an arbitrary complex in \hand{$\mathcal{M}$}$_n$(k),
we'll first have to resolve it by a sK complex, and then apply (U).

We expect \uline{that a map} $X_{*} \longrightarrow Y_{*}$ \uline{in} \hand{$\mathcal{M}$}$_n$(k) \uline{is a} \xout{\unsure{quasiis}} \uline{weak equiva-}%
\uline{lence iff the corresponding map for} LH. \uline{is a quasiisomorphism}, in other
words we expect the functor (U') to be ``conservative'': a map in the
first category is an isomorphism iff it's image in the second one is.
A second main feature we expect from linearization, is that in the case
k=\underline{Z} it corresponds to the usual abeli\unsure{an}ization of homotopy types, This
statement, when made more specific as in (w) page 454, decomposes into
two distinct ones. O\over{\unsure{n}}{n}e is of significance over an arbitrary ring k,
and states that \uline{for a} sK \uline{complex} X$_{*}$ , \uline{the inclusion} (coming from (N))
\[ \text{(V)} \qquad \mathrm{L}_{\mathrm{pt}}(X_{*}) \longrightarrow \mathrm{L}(X_{*}) \] \margin{x}
\uline{from} L$_{\mathrm{pt}}$ \uline{into it's completion}, \uline{when viewed as a map of chain complexes}
\uline{in Ab$_k$} (using the simplicial differenti\xout{\unsure{t}}al operator, or passing to the
corresponding ``normalized'' chain complexes first) \uline{is a quasi-isomorphism}.
\over{w}{W}hether this is \inserted{always} so or not, \over{not}{\unsure{or}} \xout{\unsure{xx}} \inserted{whether} noetherian conditions on k or some \xout{fin}
finiteness conditions for the components of X$_{*}$ are needed, looks like a
rather standard question of linear homological algebra ! O\over{\unsure{nx}}{n} the other
hand, using the exponential embedding (K) for sections, we get another
map \xout{of} of semisimplicial k-modules
\[ \text{(V')} \qquad k^{(X\hand{$_{*}(k)$})} \longrightarrow \mathrm{L}(X_{*}) \quad , \] \margin{x}
and here \uline{the question again} (the expectation I might say?\hand{)} \uline{is} \deleted{\xout{that}} \inserted{whether} this
\uline{is \deleted{again} \hand{a} quasiisomorphism}. This would just mean \xout{that} (if coupled with
quasi-isomorphy of (V)) that the homology invariants (U) are just the
usual homology invariants of the discrete homotopy type modelled by
X$_{*}$(k), but with coefficients not in \underline{Z} , but in k. We certainly do ex-%
pect this to be true for k = \underline{Z} - which was the content of ``question 1''

\margin{36}
